\chapter{Preface and Structural Architecture}
\label{ch:00}

\section{The Central Mathematical Structure: The Stepanov--Bost Slope Principle}
\label{sec:00:vision}

The Generalized Riemann Hypothesis (GRH) for Dirichlet $L$-functions is classically formulated as a spectral assertion about the zeros of meromorphic Euler products in the critical strip $0 < \Re(s) < 1$. Yet the unconditional proof discovered and machine-verified in Lean~4 by Levent Alp\"oge (working with Claude) in November~2026~\cite{Alpoge2026} reveals a far deeper structural phenomenon. Rather than attacking the zeros of $L(s,\chi)$ directly through contour integration in the critical strip, explicit zero-density bounds, or spectral trace formulas for $\GL_1$, the proof embeds the distribution of primes in arithmetic progressions into the Arakelov geometry of a growing tower of definite unitary period lattices over an auxiliary abelian CM field $E$, and extracts the critical-line exponent $1/2$ as an immediate shadow of a finite-dimensional slope minimax.

This monograph reconstructs the entire proof—for all Dirichlet $L$-functions $L(s,\chi)$ simultaneously, with the Riemann Hypothesis for $\zeta(s)$ as the trivial-character specialization—as a self-contained, publication-grade treatise organized around a single unifying mathematical object:
\begin{quote}
\emph{An Arakelov-filtered automorphic period tower over an abelian CM field $E$, governed by the Stepanov--Bost slope principle.}
\end{quote}
To understand why GRH emerges naturally from this structure as \cref{cor:01:grh} of our Main Structural Theorem (\cref{thm:01:main-structural}), it is helpful to recall the two classical pillars fused inside the mechanism:

\begin{enumerate}[label=(\roman*)]
\item \emph{Stepanov's polynomial method over finite fields~\cite{Stepanov1969}.} In Bombieri's streamlined form of Stepanov's proof of the Weil conjectures for curves over $\FF_q$, one constructs an auxiliary function of degree $D$ vanishing to high order $m$ at every $\FF_q$-rational point, and compares the total order of vanishing (a local divisibility lower bound) with the global degree $D$ (an upper bound on the height of the polynomial space). Balancing the number of points against the derivative order $m \asymp \sqrt{q}$ yields the square-root error term $O(\sqrt{q})$ without ever invoking \'etale cohomology.

\item \emph{Bost's slope inequality in Arakelov geometry~\cite{Bost1996}.} Over a number field $E$, the role of polynomial degree is played by the Arakelov slope $\mu(M)$ of a Hermitian $\OO_E$-lattice, while the role of high-order vanishing at rational points is played by a collection of nonzero period evaluations $\Phi_i : M_i \to N_i$ that are simultaneously \emph{divisible} at inert primes of $E$ (via height-$2$ formal group calculus and Lubin--Tate biextensions) and \emph{bounded} at the archimedean places (via the Weil representation and seesaw identities). If $\Phi$ is nonvanishing on a subspace of controlled codimension, Bost's slope inequality forces
\[
  \mu_{\min}(M) \;\le\; \mu_{\max}(N) + \log \|\Phi\|,
\]
up to controlled frame-codimension and normalization terms.
\end{enumerate}

In the setting of this book, we construct a tower of Hermitian $\OO_E$-lattices of alternating even dimensions $N_j = k + j(t+1)$ ($0 \le j \le m$) carrying definite unitary theta lifts, together with an evaluation pairing into a dual tower of automorphic row lattices of rank $D_{\mathrm{all}}$. Crucially, the period matrix $R$ of the bottom automorphic forms sits simultaneously as the \emph{source} of a left-hand covolume upper bound and as the \emph{target} of a right-hand near-orthogonality lower bound. Placing these two estimates back-to-back yields what we call the \emph{Two-Sided Arakelov Slope Sandwich}: the unknown archimedean period term $\mu(R R^\vee)$ cancels \emph{identically} between the two faces of the sandwich! What survives is a pure arithmetic inequality relating:
\begin{itemize}
\item on the \emph{left-hand side}, the \emph{inert-prime divisibility functional} $\ASK(x, m)$ over inert primes $p$ of $E$ below the scale $a = x b$, weighted by an additive Behrend-thick family of CM places; and
\item on the \emph{right-hand side}, the \emph{analytic main term} $\MAIN(x, m)$ (arising from the Siegel--Weil mass formula, the Chowla--Selberg period formula, and the degree of unitary Shimura flag towers) plus four structural error budgets of sizes $x/m$, $x m / k$, $m$, and $k / n_0^2$ (up to factors of $x^{O(\varepsilon)}$).
\end{itemize}
By the \emph{Balance Identity} (\cref{thm:03:balance-identity}), the difference $\ASK - \MAIN$ expands directly as a positively weighted average of centered prime-power ramp sums
\[
  G(\chi, y) \;=\; \sum_{n \le y} \Lambda(n)\chi(n)\!\left(\frac{y}{n}-1\right) + \text{principal term}
\]
for $\chi$ ranging over the cyclic subgroup $\langle\psi\rangle$ generated by a fixed primitive Dirichlet character $\psi$. Minimizing the four competing error scales
\[
  \frac{x}{m} \;+\; \frac{x m}{k} \;+\; m \;+\; \frac{k}{n_0^2}
\]
forces the unique geometric optimum
\[
  m \;\asymp\; x^{1/2 - \varepsilon/16}, \qquad n_0 \;\asymp\; x^{2 + 8\varepsilon}, \qquad k \;\asymp\; x^{9/2 + 10\varepsilon},
\]
at which the competing terms $x/m$ and $m$ equilibrate at $2\sqrt{x}$! Thus the pointwise Main Inequality $(\star)$ (\cref{thm:06:main-ineq-star}) immediately produces the one-sided ramp-sum bound $G_\psi(x) \le x^{1/2 + 9\varepsilon}$ across arbitrarily large exponential windows $[Y, Y^{3/2}]$. Finally, a purely elementary \emph{Window Landau Theorem} (\cref{thm:02:window-landau}) converts this one-sided bound on $[Y, Y^{3/2}]$ into the zero-free half-plane $\Re(s) > \frac{1}{2} + 9\varepsilon$ for every $\chi \in \langle\psi\rangle$. Letting $\varepsilon \to 0^+$ and applying the functional equation yields GRH.

In this light, the exponent $1/2$ in the Generalized Riemann Hypothesis is neither an accident of analytic continuation nor an unmotivated artifact of the Weil conjectures for curves: it is the unique minimax of the Arakelov-filtered period tower length $m \asymp \sqrt{x}$ balancing the derivative-order truncation error $x/m$ against the tower-height growth $m$.

\section{Relationship to the Machine-Verified Archive and Mathematical Optimizations}
\label{sec:00:provenance}

The original Lean~4.33.1 formalization~\cite{Alpoge2026} (comprising $175{,}187$ files and $\approx 17.37$ million lines of code, verified with zero unproven placeholders and zero custom axioms beyond standard constructive-classical Lean foundations) and prior human-readable guides~\cite{Rubinstein2026,Jaswal2026} necessarily reflect the engineering constraints of formal theorem proving. In particular, the Lean source tree contains extensive formalization scaffolding (intermediate stand-in interfaces, legacy alternative record bundles, and manually weakened surrogates for deep published theorems in automorphic representation theory and $p$-adic Hodge theory).

The goal of this $\approx 300$-page monograph is twofold:
\begin{enumerate}[label=(\arabic*)]
\item \emph{Architectural clarity and true mathematical optimization.} We strip away all formalization scaffolding and reorganize the thirty-six chapters of the archive-mirroring exposition into twenty tightly integrated chapters grouped into six logical parts. Whenever a major classical theorem in the published literature (such as the construction of polarized automorphic Galois representations~\cite{BLGGT2014,Shin2011,ChenevierHarris2013,Caraiani2012,Caraiani2014}, the local Langlands correspondence for $\GL_n$~\cite{HarrisTaylor2001,Henniart2000,BushnellHenniart2005,BushnellHenniart2011}, Kato's Euler-system vanishing theorem~\cite{Kato2004}, or Jiang--Sun--Zhu's archimedean Bessel multiplicity-one theorem~\cite{JiangSunZhu2010}) directly supplies an interface that the Lean archive was forced to re-engineer in artificially weakened form, we cite the published theorem precisely via \texttt{[Literature]} and verify every single hypothesis explicitly in the text.

\item \emph{Zero-hallucination rigor.} Every newly invented mathematical simplification that departs from the Lean proof route is explicitly flagged with \texttt{[New Simplification]}, accompanied by a $100\%$ self-contained, gap-free proof, and cataloged in \cref{app:A} as well as the companion optimization log. Conversely, whenever a tempting conceptual shortcut fails under quantitative stress-testing—most notably the attempted elimination of the Behrend thick set $T_2$, the attempted replacement of the explicit cohomology bound for $\GL_N(\OO_L)$ by Postnikov/motivic arguments, the attempted replacement of rank-$4$ unitary Hida propagation by parallel-weight modularity lifting, and the attempted replacement of finite-level adjoint rigidity by big-image theorems—we honestly retain the verified constructive argument, explain the exact mathematical obstruction, and link every theorem to its exact verified declaration in the Lean~4 archive via \texttt{[Lean Archive]} tags.
\end{enumerate}

The nine principal mathematical optimizations proved and verified in this book are summarized below and audited systematically in \cref{app:A}:
\begin{enumerate}[label=\textbf{O\arabic*.}]
\item \emph{Classical Automorphic--Galois Interface} (\cref{thm:01:classical-automorphic-galois}): Replaces the archive's artificially weakened Galois-representation and local-parameter constructions with published theorems of Harris--Taylor, Henniart, Bushnell--Henniart, Labesse, Moeglin--Waldspurger, BLGGT, Shin, Chenevier--Harris, Caraiani, and Bella\"iche--Chenevier, with explicit proofs of all four compatibility subtleties (Speh cyclotomic shifts, half-integral twists and polarization signs, crystallinity at $\ell \in \{2,3\}$ and ramified $2$-adic places, and level-zero inertia matching).
\item \emph{Single-Copy ($k=1$) Landherr Isometry and Elementary Cayley Transform} (\cref{thm:07:landherr-single-copy}): Eliminates fourfold tensoring of Hermitian spaces and fourfold orthonormal frame construction in the proof of inert-prime theta integrality, directly establishing the Weil transfer scalar $c^8 = +\det(u)^{+4}$ via Landherr's theorem~\cite{Landherr1936} and Hilbert reciprocity.
\item \emph{Serre--Tate and Lubin--Tate/Honda CM Good Models} (\cref{thm:07:cm-good-models}): Compresses the multi-stage good-model construction at inert primes $p \ge 5$ to a single transparent argument via centralizer rigidity $\operatorname{Cent}_{\End(\widetilde{\mathcal{E}})}(\OO_K) = \OO_K$ and quadratic twisting.
\item \emph{Equivariant Bicoboundary Vanishing, Proper-Cone Euclidean-Projection LP Duality, and Sharpened Inert Constant $8$} (\cref{thm:07:biextension-rigidity,thm:08:lp-duality,thm:08:inert-deriv}): Replaces Hahn--Banach and Tychonoff compactness in the cyclic orbital linear-programming engine with a self-contained $14$-line Bolzano--Weierstrass proper-cone Euclidean projection argument, proves equivariant biextension rigidity in six lines, and provides an eight-line one-shot dual test function sharpening the odd-level inert derivative constant from $11$ to $8$.
\item \emph{Basis-Free Hilbert--Schmidt Paley--Zygmund Frame Codimension Inequality and Direct Sharpened Bost Slope} (\cref{thm:05:frame-codim,thm:05:bost-slope}): Establishes the strictly tighter codimension bound $\dim(W)/D \le \frac{\eta}{1+\eta} + 4\varepsilon + 2\sqrt{e}$ in ten basis-free lines, removes the legacy $-\log|W(u_0)|$ loss directly from the multi-map Bost inequality (rendering the global error $\mathrm{Err}_5$ independent of the translate $u_0$), and proves isometric column compression once for the entire book.
\item \emph{Unified Master Absorption Lemma and Four-Class Loss Consolidation} (\cref{thm:10:four-class-losses}) together with \emph{Four-Block Arithmetic Size Consolidation} (\cref{thm:12:size-inputs}): Merges all eleven non-inert loss sums into four clean prime-class regimes with locked exponent $e_L = 0$ and consolidates the eleven arithmetic size inputs into four natural geometric blocks.
\item \emph{Covolume--RTF Two-Sided Duality, Bessel Multiplicity One via Jiang--Sun--Zhu, Six-Cluster Classical Geometric Box, and Lossless Four-Theorem Dissolution of the Definite-Side Interface} (\cref{thm:13:volume-pq,thm:14:covolume,thm:15:diag-wave-grid,thm:16:offdiag-near-orth,thm:17:definite-assembly}): Re-homes the diagonal rotated wave-grid bound directly after four-way Poincar\'e unfolding, replaces the ten-page archimedean uniqueness argument with Jiang--Sun--Zhu~\cite{JiangSunZhu2010} plus a $1.5$-page local scalar rigidity proof, organizes all $42$ classical geometric inputs into six clean thematic clusters, and dissolves the $46$-component definite-side hypothesis record losslessly into four conceptual theorems.
\item \emph{Four-Step Compressed Cohomological Bound for $\GL_N(\OO_L)$ via Borel and van der Kallen} (\cref{thm:19:lemma-k}): Proves the primitive odd-degree rank bound ($\le r_2(L)$) in six lines using van der Kallen homological stability~\cite{vanderKallen1980}, Borel's real cohomology ring~\cite{Borel1974}, and the exact $\ell^1$ eigenvalue of $\ell$-fold block replication on single-trace invariant forms (giving $\ell^r$ on $r$-fold cup products).
\item \emph{Integrated Selmer Filtration, Unified Four-Step Graded Deformation Machine, and One-Shot Kato Vanishing over $K/\QQ$} (\cref{thm:20:graded-length,thm:20:vwin-complete}): Unifies the parity splitting, four deformation-counting strata, and single-pass $\sqrt{b} = \exp(\frac{1}{2}x^\varepsilon)$ assembly with zero cross-chapter duplication, and collapses the two-weight Euler-system vanishing argument to four lines via Kato~\cite[Theorem~14.2(1)]{Kato2004} and Jacquet--Shalika~\cite{JacquetShalika1976} over $K/\QQ$.
\end{enumerate}

\section{Global Dependency Graph of the Monograph}
\label{sec:00:dag}

The twenty chapters are organized into six parts whose directed acyclic dependency structure is completely transparent:
\begin{center}
\renewcommand{\arraystretch}{1.18}
\begin{tabular}{@{}lll@{}}
\toprule
\textbf{Part} & \textbf{Chapters} & \textbf{Core Structural Role} \\
\midrule
\textbf{I} & Ch.~\ref{ch:01}--\ref{ch:04} & Main Structural Theorem, GRH as \cref{cor:01:grh}, Window Landau, Balance Identity, Host Supply \\
\textbf{II} & Ch.~\ref{ch:05}--\ref{ch:06} & Squared-Norm Bost Slope, Tight Frame Engine, Two-Sided Sandwich, Main Inequality $(\star)$ \\
\textbf{III} & Ch.~\ref{ch:07}--\ref{ch:10} & Period Towers, Height-$2$ Formal Groups, Inert Derivatives \& LP, Glued Lattices, $11$ Losses \\
\textbf{IV} & Ch.~\ref{ch:11}--\ref{ch:14} & Theta Proportionality, Seesaw Periods, Analytic Size, Signature-$(p,q)$ Volume, Covolume Descent \\
\textbf{V} & Ch.~\ref{ch:15}--\ref{ch:17} & Poincar\'e Unfolding, Wave Grid, Off-Diagonal Orthogonality, Limit Multiplicity, $\ell\in\{2,3\}$ \& Green \\
\textbf{VI} & Ch.~\ref{ch:18}--\ref{ch:20} & Twisted-Place Integrality, Selmer Filtration, Minus Part, Cohomological Bound, Graded Deformations \\
\bottomrule
\end{tabular}
\end{center}
The global flow is strictly acyclic:
\begin{itemize}
\item \textbf{Part~VI} (Ch.~\ref{ch:18}--\ref{ch:20}) uses the Classical Automorphic--Galois Interface (\cref{thm:01:classical-automorphic-galois}) to establish the unconditional window bound $\le \exp(\frac{1}{2}x^\varepsilon)$ on the higher Tate--Shafarevich lengths (\cref{thm:20:vwin-complete}) and converts it into the twisted-place integrality bound (\cref{thm:18:twisted-integrality}).
\item \textbf{Part~III} (Ch.~\ref{ch:07}--\ref{ch:10}) supplies the \emph{local-to-global divisibility} half: inert-prime theta integrality, height-$2$ derivative divisibility, glued Hermitian lattices, and the four-class loss bound $e_L = 0$.
\item \textbf{Part~IV} (Ch.~\ref{ch:11}--\ref{ch:14}) supplies the \emph{upper slope} half of the Arakelov sandwich: theta proportionality, seesaw period transport, analytic kernel bounds, mean-height lower bounds on signature-$(p,q)$ flag towers, and Arakelov covolume descent on deep Shimura covers.
\item \textbf{Part~V} (Ch.~\ref{ch:15}--\ref{ch:17}) supplies the \emph{lower slope / spectral} half of the Arakelov sandwich: four-way Poincar\'e unfolding, rotated wave-grid diagonal bounds, off-diagonal near-orthogonality, quantitative limit multiplicity, and definite-side closure at all primes including $\ell \in \{2,3\}$.
\item \textbf{Part~II} (Ch.~\ref{ch:05}--\ref{ch:06}) fuses Parts~III--VI inside the Two-Sided Arakelov Slope Sandwich to prove the pointwise Main Inequality $(\star)$ (\cref{thm:06:main-ineq-star}) and close the global tower family (\cref{thm:06:global-assembly}).
\item \textbf{Part~I} (Ch.~\ref{ch:01}--\ref{ch:04}) converts $(\star)$ via the Balance Identity, host CM supply, and Window Landau Theorem into the Main Structural Theorem (\cref{thm:01:main-structural}) and \cref{cor:01:grh} (GRH).
\end{itemize}

\section{Locked Global Parameters, Normalizations, and Notation}
\label{sec:00:notation}

To prevent any normalization drift across the twenty chapters, all global parameters, slope conventions, and Frobenius normalizations are locked in \cref{tab:00:constants} and enforced uniformly throughout the monograph.

\begin{table}[ht]
\centering
\caption{Locked global scales, exponents, and arithmetic normalizations.}
\label{tab:00:constants}
\renewcommand{\arraystretch}{1.15}
\small
\begin{tabular}{@{}lll@{}}
\toprule
\textbf{Symbol} & \textbf{Locked Definition / Range} & \textbf{Origin} \\
\midrule
$\varepsilon$ & $0 < \varepsilon < 1/52$ (from $200\varepsilon + \frac{2}{13} \le 4$) & Ch.~\ref{ch:04}, \cref{thm:04:supply} \\
$\lambda$ & $\lambda = 3/2$ (contracted from $[W,W^2]$ at cutoff $401$ to $[Y,Y^{3/2}]$ at $400$) & Ch.~\ref{ch:02} \& Ch.~\ref{ch:04} \\
$L,\, b,\, \sqrt{b}$ & $L = x^\varepsilon$, $b = \exp(x^\varepsilon)$, $\sqrt{b} = \exp(\frac{1}{2}x^\varepsilon)$ (never $\exp(x^{\varepsilon/2})$!) & Ch.~\ref{ch:03}, \ref{ch:06}, \ref{ch:20} \\
$a,\, X',\, b'$ & $a = x b$, $X' = \exp(x^{\varepsilon/3})$, $b' = a/X'$ & Ch.~\ref{ch:03} \& Ch.~\ref{ch:06} \\
$\delta_1,\, \delta_2$ & $\delta_1 = x^{-\varepsilon/8}$ (top density), $\delta_2 = |T_2|/n_0 \ge x^{-\varepsilon}$ (Behrend density) & Ch.~\ref{ch:03} \& Ch.~\ref{ch:04} \\
$n_0,\, f$ & $n_0 = |A_0| = 2^e \in [x^{2+8\varepsilon}, 2x^{2+8\varepsilon}]$, $f = 2 n_0$ & Ch.~\ref{ch:03} \& Ch.~\ref{ch:06} \\
$m,\, k$ & $m \in [x^{1/2-\varepsilon/16}, 2x^{1/2-\varepsilon/16}]$ (even), $k \in [x^{9/2+10\varepsilon}, 2x^{9/2+10\varepsilon}]$ (even) & Ch.~\ref{ch:03} \& Ch.~\ref{ch:06} \\
$t,\, N_j$ & $t$ odd ($t \le x^{400}$), $N_j = k + j(t+1)$ (all even, $0 \le j \le m$) & Ch.~\ref{ch:06} \& Ch.~\ref{ch:07} \\
$C_{\mathrm{abs}}, C_0, e_D, e_L$ & $(C_{\mathrm{abs}}, C_0, e_D, e_L) = (2, 3, 2, 0)$; $\mathrm{CErr} = e_D + e_L + 3 = 5$ & Ch.~\ref{ch:06} \& Ch.~\ref{ch:10} \\
Exponent ledger & $2 \to 5 \to \frac{53}{8} \to \frac{61}{8} \to 8 \to 9$, yielding $G_\psi(x) \le x^{1/2 + 9\varepsilon}$ & Ch.~\ref{ch:03} \& Ch.~\ref{ch:06} \\
\bottomrule
\end{tabular}
\end{table}

\paragraph{Slope normalization.}
For a Hermitian $\OO_L$-lattice $M$ of rank $d \ge 1$ over a totally complex number field $L$ of degree $2 r_2(L) = [L:\QQ]$ and discriminant $\Delta_L$, we use the \emph{squared-norm Arakelov slope} throughout Parts~I--III and Ch.~\ref{ch:14}--\ref{ch:20}:
\begin{equation}
\label{eq:00:squared-slope}
  \mu(M) \;=\; 2\,\widehat{\mu}_{\mathrm{Bost}}(M) - \log 2 + \frac{\log|\Delta_L|}{[L:\QQ]}, \qquad
  \operatorname{lcv}(X, M) \;=\; -d\,\mu(M), \qquad \mu(M^\vee) \;=\; -\mu(M).
\end{equation}
In \cref{ch:13} alone, the intersection-theoretic volume theorem on unitary Shimura flag bundles naturally uses the \emph{unsquared} Arakelov mean height $\widehat{\mu}_{\mathcal{T}}$; the exact factor-of-$2$ bridge converting $\widehat{\mu}_{\mathcal{T}}$ into $\mu$ is executed explicitly at four marked checkpoints in \cref{ch:14}.

\paragraph{Frobenius, Satake, and Hodge--Tate normalizations.}
We adhere uniformly to the following conventions:
\begin{itemize}
\item Arithmetic Frobenius has degree $+1$ ($\deg(\Frob_w) = +1$); under the local Artin isomorphism $\rec_w$, \emph{geometric} Frobenius $\Frob_w^{-1}$ corresponds to a uniformizer $\varpi_w$.
\item Unramified Satake parameters $(\alpha_1,\dots,\alpha_a)$ of an irreducible representation $\pi_w$ of $\GL_a(E_w)$ are normalized unitarily so that the geometric Frobenius eigenvalues on the associated Weil--Deligne representation are $q_w^{(a-1)/2}\alpha_i$ (the unitary Tate-shift normalization of \cite{BLGGT2014,HarrisTaylor2001}).
\item The $\ell$-adic cyclotomic character $\varepsilon_\ell$ has Hodge--Tate weight $\HT(\varepsilon_\ell) = \{-1\}$, matching Deligne's cohomological convention for $H^1_{\text{\'et}}$ of elliptic curves (weights $\{0, -1\}$) and modular forms of weight $k \ge 2$ (weights $\{0, 1-k\}$).
\end{itemize}

\section{Analytic Rigor Manifesto: Six Classical Pitfalls Explicitly Guarded Against}
\label{sec:00:manifesto}

History shows that attempted proofs of the Riemann Hypothesis routinely founder on subtle analytic slippages—sign inversions, uncontrolled contour tails, illegal $L^1$-to-$L^\infty$ passages, or silent quantifier collapses. A systematic forensic audit of prior erroneous analytic manuscripts in the companion repository identified six recurring failure modes. Every analytic argument in this book has been explicitly stress-tested against all six:

\begin{enumerate}[label=\textbf{G\arabic*.}]
\item \emph{No $L^1 \Rightarrow L^\infty$ leap and no uncontrolled exterior tails.}
In \cref{ch:02}, the Window Landau inequality flows strictly from \emph{pointwise} window upper bounds $g(u) \le R(u)$ to an $L^1$ integral bound, while the exterior tail $(0,\infty) \setminus [U,V]$ is controlled by an exact two-sided Chernoff bound requiring neither Stirling's formula nor asymptotic approximations. In \cref{ch:04}, passage from mean-square ($L^2$) prime-sum bounds to uniform ($L^\infty$) window bounds uses the exact right-derivative Sobolev--Cauchy--Schwarz self-similarity identity (\cref{lem:04:sobolev-selfsim}), never an unverified $L^1$-to-$L^\infty$ inference.

\item \emph{No discontinuous test kernels and no global critical-strip contour shifts.}
The ramp kernel $(y/n - 1)_+$ is continuous with piecewise-smooth Mellin transform $(s(s-1))^{-1}$, explicitly desingularized at $s = 1$ via divided differences. In \cref{ch:02}, singularity detection uses only \emph{local} Cauchy disks $\overline{B}(z, r') \subset \{\Re(s) > 0\}$; in \cref{ch:04}, Perron shifts stay strictly inside flat zero-free rectangles $\{\Re(s) > 1 - 1/400,\, |\Im(s)| \le Q\}$, preventing corner leakage.

\item \emph{Signed identities locked and independently cross-checked.}
Every critical sign—the strictly positive real-pole residue $+\frac{S(\beta)}{\beta(1-\beta)} > 0$ of the symmetrized Laplace transform $\Phi$, the $+a\ell(\chi_C)$ term of the Balance Identity, the $\iota_d$-thick preimage sign $+1$ preventing odd-character sign flips, the Weil transfer identity $c^8 = +\det(u)^{+4}$ in \cref{thm:07:landherr-single-copy}, and the factor-$2$ slope conversion $\mu = 2\widehat{\mu}_{\mathrm{Bost}}$—is fixed in canonical coordinates and cross-checked on explicit test cases.

\item \emph{No critical-strip Euler products and no unbounded biorthogonal spectral bases.}
Nowhere in the proof is an Euler product manipulated inside the critical strip, nor are zero-spacing assumptions or non-Schauder spectral expansions invoked. Every step operates on finite-rank Hermitian lattices and finite sums truncated at explicit scales.

\item \emph{Strict quantifier hierarchy ($\varepsilon$-separation).}
All structural exponents ($e_D = 2$, $e_L = 0$, $\mathrm{CErr} = 5$, and the final ramp exponent $9$) are proved uniformly \emph{before} $\varepsilon \in (0, 1/52)$ is fixed, while all asymptotic thresholds and multiplicative constants explicitly carry their dependence on $(\varepsilon, \psi)$.

\item \emph{Complete boundary-regime splits and strict admissibility.}
Every one-sided Laplace/Mellin integral is supported strictly on $(0,\infty)$ with explicit $L^1$ domination; every character family enforces the admissibility condition ($N_i = 1$ or $\chi_i \ne 1$); and every piecewise prime-power estimate explicitly splits boundary regimes (including $p \le y$ versus $p > y$ in \cref{ch:04} and the $D = 0$ guard in \cref{ch:05}).
\end{enumerate}

With these conventions and safeguards in place, we turn immediately in \cref{ch:01} to the Main Structural Theorem, the emergence of GRH as \cref{cor:01:grh}, and the Classical Automorphic--Galois Interface.
